\subsection{Strictness of the first-valuation bound in odd characteristic}
\label{subsec:odd-positive-excess-strict}

For odd characteristic the first-valuation estimate
\eqref{eq:positive-square-first-valuation} is strict.
The proof keeps the equality cases in both the genus identity and the
resolution-loss estimate. These equality cases would force cusps of
type $(2,p)$ on the normalized root curve. A regular logarithmic form
on the normalized surface then has a conormal leading coefficient
that vanishes above each cusp, whereas the zero self-intersection
of the smooth strict transform requires that coefficient to be
nowhere zero.

\begin{theorem}\label{thm:odd-positive-excess-strict}
Let $X$ be a normal integral projective surface over an algebraically
closed field $k$ of odd characteristic $p$, and let $C\subset X$ be
a smooth integral effective Cartier divisor contained in the smooth
locus, with $d=C^2>0$.
Let $M\ne\OO_X$ have order $p$ and satisfy
$M|_C\simeq\OO_C$.
Suppose that its logarithmic form has exact order
$m=pa$, where $p\nmid a$.
If $[M]\in p\operatorname{NS}(X)$, then
\[
 \delta:=2g(C)-2-pa\,d>\frac{p\,d}{2(p-1)}.
\]
Equivalently,
\[
 \delta\geq
 \left\lfloor\frac{p\,d}{2(p-1)}\right\rfloor+1.
\]
In particular the strict inequality holds for every algebraically
trivial such $M$.
\end{theorem}

\begin{proof}
Theorem~\ref{thm:positive-excess-picard} gives the weak inequality.
We shall exclude
\begin{equation}\label{eq:odd-endpoint-assumption}
 2(p-1)\delta=pd.
\end{equation}

\emph{Step 1: the root cover and equality in the losses.}
We first choose the root on the original surface.
By the Picard-variety argument in the proof of
Theorem~\ref{thm:positive-excess-picard}, the reduced identity
component of the Picard scheme is an abelian variety, and
multiplication by $p$ is surjective on its $k$-points.
Choose $L_0$ representing a N\'eron--Severi root of $M$.
Then $M\otimes L_0^{-p}$ is algebraically trivial, so it has
an algebraically trivial $p$th root $L_1$.
Thus
\[
 A:=L_0\otimes L_1,\qquad A^{\otimes p}\simeq M.
\]
Resolve $X$ projectively by a morphism $\rho:S\to X$ that is an
isomorphism near $C$.
The curve, its self-intersection, its genus and the logarithmic
order are unchanged.
Normality gives $\rho_*\OO_S=\OO_X$, and projection formula makes
pullback on Picard groups injective.
We retain the notation $C,M,A$ for their pullbacks to the smooth
surface $S$.

Let $\pi:N\to S$ be the normalization of the $\mu_p$-torsor of $M$,
and put $D=\pi^*C$ scheme-theoretically.
Lemma~\ref{lem:positive-square-root-genus} says that $N$ is normal,
integral and projective, that $D$ is a reduced integral Cartier
divisor, and that
\[
 D^2=pd,\qquad g(\overline D)=g(C),\qquad
 (p-1)\delta=\ell+2\Delta,
\]
where
\[
 \Delta=p_a(D)-g(C),\qquad \ell\geq0.
\]
The normalization map $\phi:\overline D\to C$ is the degree-$p$
purely inseparable map with function field $k(C)^{1/p}$.

Use the resolution $f:T\to N$ constructed in the proof of
Lemma~\ref{lem:positive-square-resolution-loss}, and write $F$
for the smooth strict transform of $D$.
By \eqref{eq:odd-endpoint-assumption},
\[
 F^2\geq D^2-4\Delta
       \geq pd-2(p-1)\delta=0.
\]
If $F^2>0$, Lemma~\ref{lem:positive-square-resolved-transfer},
applied to the chosen root $A$, gives exact logarithmic order $a$
and forces $p\mid a$.
Therefore
\begin{equation}\label{eq:odd-endpoint-zero-square}
 F^2=0.
\end{equation}
It follows that equality holds throughout
\[
 pd=D^2-F^2\leq4\Delta\leq2(p-1)\delta=pd.
\]
In particular,
\begin{equation}\label{eq:odd-endpoint-genus}
 \Delta=\frac{pd}{4}>0,\qquad \ell=0.
\end{equation}

We retain the equality information in the resolution construction.
Its first stage is the resolution $f_0:T_0\to N$ with no exceptional
smooth rational curve of square $-1$.
Write
\[
 f_0^*D=D_0+E,\qquad
 \lambda_0=D^2-D_0^2=D_0\cdot E=-E^2,\qquad
 \Delta_0=p_a(D_0)-g(C).
\]
The proof of Lemma~\ref{lem:positive-square-resolution-loss} gives
\[
 2(\Delta-\Delta_0)
   =\lambda_0+K_{T_0}\cdot E\geq\lambda_0\geq0.
\]
In its second stage, let $\mu_i\geq2$ be the multiplicities
at the successive singular points of the strict curve.
Then
\[
 \sum_i\frac{\mu_i(\mu_i-1)}2=\Delta_0,\qquad
 D^2-F^2=\lambda_0+\sum_i\mu_i^2.
\]
Consequently
\begin{align*}
 pd
 &=\lambda_0+\sum_i\mu_i^2\\
 &\leq2(\Delta-\Delta_0)+4\Delta_0\\
 &=4\Delta-2(\Delta-\Delta_0)\\
 &\leq4\Delta=pd.
\end{align*}
Equality of the endpoints forces
\begin{equation}\label{eq:odd-endpoint-resolution-equalities}
 \Delta=\Delta_0,\qquad \lambda_0=0,\qquad
 \mu_i=2\ \text{for every }i.
\end{equation}

Since $E^2=0$, we have $E=0$.
Indeed, a nonzero effective exceptional divisor is orthogonal to
the pullback of an ample divisor on $N$ and has negative square
by the Hodge index theorem.
It cannot be numerically zero, since its intersection with an
ample divisor on $T_0$ is positive.
Thus $D$ avoids the singular locus of $N$:
if a singular point belonged to $D$, a local equation of $D$
would vanish along every exceptional curve above it and give
a positive coefficient in $E$.
A resolution of a singular normal surface point has an
exceptional curve; otherwise proper quasi-finiteness and
normality would make the resolution an isomorphism there.
We conclude that $N$ is smooth on a neighborhood of $D$.
This conclusion uses $\lambda_0=0$, rather than merely $\ell=0$.

\emph{Step 2: the singularities of the root curve.}
Let $\beta$ be the first tangential coefficient of the logarithmic
form of $M$, and let $Z$ be its zero divisor on $C$.
It is a nonzero section of
$\omega_C\otimes\OO_C(-paC)$, so
\[
 \deg Z=\delta.
\]
Where $\beta$ does not vanish, the local normalization chart
in Lemma~\ref{lem:positive-square-root-genus} has $v_x$ invertible.
The chart and its inverse-image curve are smooth there.
Since $\pi$ is radicial,
\begin{equation}\label{eq:odd-endpoint-count}
 \#\operatorname{Sing}(D)
 \leq\#\operatorname{Supp}(Z)\leq\deg Z.
\end{equation}

Let $P$ be a singular point of $D$.
By \eqref{eq:odd-endpoint-resolution-equalities}, its multiplicity
in the smooth surface $N$ is two.
The normalization has a unique point above $P$, because
$\phi:\overline D\to C$ is radicial.
Completion of the finite normalization injection embeds
$\widehat{\OO}_{D,P}$ in $k[[s]]$, the completed ring of that
normalization point.
The normalization defect is unchanged by completion
\cite[Tags~0C1R and~0C3V]{Stacks}.
Thus the completed plane curve has one branch.

We spell out the multiplicity-two calculation, using $p\ne2$.
The multiplicity of a plane branch is the minimum valuation of
a general linear combination of its two plane parameters.
For this valuation formula, compare the quotient by such a
parameter with the quotient in the finite normalization:
the finite-length normalization defect contributes equal
kernel and cokernel lengths and cancels.
There is therefore an element $x$ of valuation two.
Taking the square root of its unit factor in $k[[s]]$ permits a
normalization parameter with $x=s^2$.
Since every element of the square of the maximal ideal has
valuation at least four, $x$ is part of a minimal pair of
plane parameters.
Subtract the even part of the second parameter's power series
in $s$, and divide its remaining odd part by its unit power
series in $s^2$.
These are invertible changes of complete plane coordinates.
They give a second parameter $y=s^q$ for an odd integer $q\geq3$.
An odd term must occur: otherwise the completed branch's
fraction field would be contained in $k((s^2))$.
Finite birational normalization remains birational after this
one-branch completion, since a common nonzero denominator
for its finite module remains a common denominator after
completion.
We have proved
\[
 \widehat{\OO}_{D,P}=k[[s^2,s^q]],\qquad
 D:\ y^2-x^q=0.
\]

A parameter on $C$ pulls back under $\phi$ with valuation $p$.
Thus the semigroup of $k[[s^2,s^q]]$ contains the odd integer $p$.
Its least odd value is $q$, so $q\leq p$.
The missing monomials $s,s^3,\ldots,s^{q-2}$ show that the
local normalization defect is
\[
 \Delta_P=\frac{q-1}{2}\leq\frac{p-1}{2}.
\]
The normalization exact sequence gives
$\Delta=\sum_{P\in\operatorname{Sing}(D)}\Delta_P$
\cite[Tag~0C1R]{Stacks}.
Using \eqref{eq:odd-endpoint-genus} and
\eqref{eq:odd-endpoint-count}, we obtain
\[
 \frac{p-1}{2}\deg Z
 =\Delta
 \leq\frac{p-1}{2}\#\operatorname{Sing}(D)
 \leq\frac{p-1}{2}\deg Z.
\]
Every inequality is an equality.
Hence $Z$ is reduced, there is exactly one singular point of
$D$ over each point of $Z$, and every such point has the formal
equation
\begin{equation}\label{eq:odd-endpoint-cusp}
 y^2=x^p,\qquad x=s^2,\quad y=s^p.
\end{equation}
There is at least one such cusp, since $\Delta>0$.

\emph{Step 3: a degree-zero conormal leading coefficient.}
Use frames of $A$ with transitions $g_{ij}$ and their $p$th
powers as frames of $M$.
The torsor coordinates satisfy $z_i=g_{ij}^p z_j$ and remain
regular units on $N$.
Consequently
\[
 \theta_N=\frac{dz_i}{z_i}
\]
is a global regular K\"ahler one-form on $N$, closed on its
smooth neighborhood of $D$.
Its pullback $\theta_T=f^*\theta_N$ is regular and closed on $T$.
At the generic point of $D$, the calculation in
Lemma~\ref{lem:positive-square-resolved-transfer} gives
\[
 z=b+h^a t,\qquad t^p=v,
\]
with $dh,dt$ a cotangent basis.
The derivation dual to $dt$ kills the old base ring, including
the possibly nonconstant $b$.
Thus the $dt$-coefficient of $\theta_T$ is exactly $h^a/z$.
The form is nonzero, and its exact order $j$ along $F$ satisfies
$0\leq j\leq a$.
If $j=a$, its tangential leading coefficient is nonzero.
The coefficient of order $a-1$ in its exterior derivative
then forces $p\mid a$, a contradiction.
Therefore $j<a$, and its leading restriction is conormal.

Divisibility by a power of the local equation of the smooth
curve $F$ is checked at its generic point.
The nonzero leading restriction is thus a global section
\begin{equation}\label{eq:odd-endpoint-leading-section}
 0\ne\sigma\in
 H^0\bigl(F,\OO_F(-(j+1)F)\bigr).
\end{equation}
Its tangential image vanishes generically, hence everywhere.
By \eqref{eq:odd-endpoint-zero-square}, this line bundle has
degree zero.
The section $\sigma$ must therefore be nowhere zero.

\emph{Step 4: the zero forced at each cusp.}
Fix a cusp and use the complete smooth surface coordinates
in \eqref{eq:odd-endpoint-cusp}.
Put $h_0=y^2-x^p$.
The order of $\theta_N$ along $D$ is $j$, since the resolution
is an isomorphism at the generic point of $D$.
Its coefficients are divisible by $h_0^j$ in the regular
local ring: divisibility by this power of a prime Cartier
equation is detected at the generic point, even though
the Cartier curve is singular.
After completion we may write
\[
 \theta_N=h_0^j\bigl(A(x,y)\,dx+B(x,y)\,dy\bigr)
\]
with $A,B$ regular power series.
The expression in parentheses has zero tangential image
on the normalization of $D$.
Since $x=s^2$ and $y=s^p$, that image is
\[
 2sA(s^2,s^p)\,ds.
\]
Here $2s\ne0$, so $A(s^2,s^p)=0$.
Thus $A=h_0A_1$ and
\begin{equation}\label{eq:odd-endpoint-local-form}
 \theta_N
   =h_0^j\bigl(h_0A_1(x,y)\,dx+B(x,y)\,dy\bigr).
\end{equation}

Write $r=(p-1)/2$.
The successive curve blowups at this cusp have equations
$t_i^2=x^{p-2i}$.
There are exactly $r$ of them before the strict curve is smooth.
The final chart is
\[
 y=x^r t,\qquad F:\ t^2=x.
\]
Blowups at other cusps do not change this chart.
Set $H=x-t^2$, a local equation of $F$.
Then
\[
 x=t^2+H,\qquad y=x^r t,\qquad h_0=-x^{p-1}H.
\]
Along $F$,
\[
 dx|_F=2t\,dt+dH
\]
and
\begin{align*}
 dy|_F
 &=r t^{2r-1}\,dH+(2r+1)t^{2r}\,dt\\
 &=r t^{p-2}\,dH,
\end{align*}
because $2r+1=p=0$ in $k$.
The scalar $r$ is nonzero in $k$.
Pull \eqref{eq:odd-endpoint-local-form} to this chart, divide
by $H^j$, and restrict to $F$.
The term with $h_0A_1\,dx$ still has a factor $H$ and vanishes.
The remaining conormal coefficient is
\[
 (-1)^j r\,t^{\,2j(p-1)+p-2}B(t^2,t^p).
\]
It vanishes at $t=0$, since $B$ is regular and
\[
 2j(p-1)+p-2\geq p-2>0.
\]
Multiplying a local equation by a unit does not change
this vanishing.
Faithful flatness of completion shows that it is a zero
of the original section
\eqref{eq:odd-endpoint-leading-section}.
This contradicts its nowhere-vanishing property.
Therefore \eqref{eq:odd-endpoint-assumption} is impossible,
and the weak inequality is strict.
\end{proof}

The proof uses a regular form on the normalized surface and its
behavior at the resolved cusps.
It requires no inference of semiampleness or of a fibration
from the zero self-intersection of the strict transform.
